\documentclass{article}
\usepackage{graphicx} % Required for inserting images

\setlength{\headheight}{13.5pt}
\addtolength{\voffset}{-1.5cm}
\addtolength{\hoffset}{-1cm}
\addtolength{\textwidth}{2cm}
\addtolength{\textheight}{3cm}

%Header-Footer
\usepackage{fancyhdr}
%Header Info
\lhead{Profesor: C\'esar Hern\'andez Cruz}
\rhead{\'Algebra Lineal I 2024-1}

\title{BDE: Práctica   4 Modelo Entidad - Relación}
\author{Carlos Fabián Esteva Gallegos, Andrés López González} 

\begin{document}

\maketitle

\textbf{Objetivo:} Realizar consultas en SQL a las bases de datos de los sistemas de Reclusorios y Pacientes en base de datos Postgres.

\section*{Actividad 1:}
En las bases de datos Reclusorios\_XXX y Pacientes\_XXX QUE CREASTE EN LA PRACTICA ANTERIOR, cargar 10 registros en las tablas correspondientes a las entidades fuertes y los registros que correspondan en las entidades débiles conforme a las siguientes restricciones:

a) Las llaves primarias no deben ser numéricas, sino alfanuméricas empezando por tres letras y tres números.

\begin{table}[h]
    \centering
    \begin{tabular}{|l|l|}
        \hline
        Id paciente & PAC001 \\
        \hline
        Id juez & JUE004 \\
        \hline
        Id enfermedad & ENF003 \\
        \hline
    \end{tabular}
\end{table}

b) Recuerde las reglas de integridad de dominio, de valor y referencias vistas en clase y respétalas al cargar la información.



\section*{Actividad 2: Verifique los datos capturados}

a) Consulte todos los registros de todas las tablas de Reclusorios\_XXX.


b) Consulte todos los registros de todas las tablas y Pacientes\_XXX.

\section*{Actividad 3: Realizar las siguientes consultas al Sistema Reclusorios\_XXX}

\begin{enumerate}
    \item[i.] Dado el nombre del recluso, obtener quien fue el juez que lo juzgo y cuántos días deberá cumplir por condena.
    
    \item[ii.] Dado el nombre del juez, obtener el nombre de los reclusos que ha juzgado, así como el promedio de días que suele poner por condena a estos reclusos.
    
    \item[iii.] Dado el nombre del delito, obtener el número de reclusos que están cumpliendo condena por ese mismo delito.
    
    \item[iv.] Dado el número de días obtener el nombre de los reclusos que cumplen una pena del mismo número de días.
\end{enumerate}


\section*{Actividad 4: Realizar las siguientes consultas al Sistema Pacientes\_XXX}

\begin{enumerate}
    \item[i.] Dado el nombre de un paciente, conocer el nombre de todas las enfermedades que ha padecido.
    
    \item[ii.] Dado el nombre de una enfermedad, conocer el nombre y domicilio de todos los pacientes que la han padecido.
    
    \item[iii.] Dado el nombre de una enfermedad, conocer el nombre de todos los pacientes para los cuales su padre o su madre hayan padecido dicha enfermedad.
    
    \item[iv.] Dado el nombre de un paciente, conocer todas las enfermedades que hayan padecido los hermanos y primos de dicho paciente.
    
    \item[v.] Conocer el nombre de todas las enfermedades contagiosas, y para cada una de éstas, conocer el nombre de todos los pacientes que la hayan padecido.
\end{enumerate}

\begin{verbatim}
SELECT * FROM a
INNER JOIN b ON a.key = b.key

SELECT * FROM a
LEFT JOIN b ON a.key = b.key

SELECT * FROM a
WHERE b.key IS NULL

SELECT * FROM a
WHERE a.key IS NULL

SELECT * FROM a
FULL JOIN b ON a.key = b.key

SELECT * FROM a
RIGHT JOIN b ON a.key = b.key

SELECT * FROM a
WHERE a.key IS NULL
\end{verbatim}

a) Cree las tablas y capture los datos EN LA BASE DE DATOS QUE ACABAS DE CREAR que se indican a continuación:

\begin{table}[H]
    \centering
    \begin{tabular}{|l|l|l|l|l|}
        \hline
        \textbf{id\_Empleado} & \textbf{Nombre\_Empleado} & \textbf{id\_Proy} & \textbf{id\_Proy} & \textbf{Nombre\_Proyecto} \\
        \hline
        EMP001 & Juan Pérez & PROY001 & PROY001 & Inteligencia Negocio \\
        \hline
        EMP002 & Pedro López & NULL & PROY003 & Big Data \\
        \hline
        EMP004 & Armando Contreras & NULL & PROY005 & Aseguramiento Calidad \\
        \hline
        EMP005 & Elisa Ramírez & PROY003 & & \\
        \hline
        EMP007 & Laura Rodríguez & PROY003 & & \\
        \hline
        EMP009 & Esteban Gutiérrez & PROY001 & & \\
        \hline
        EMP010 & Claudia Méndez & PROY003 & & \\
        \hline
    \end{tabular}
\end{table}

b) Investigue la sintaxis para diferentes tipos de join y resuelva las siguientes consultas:

\begin{enumerate}
    \item[i.] Obtener todos los Empleados que tienen Proyectos asignados
    
    \item[ii.] De todos los Empleados, quienes tienen Proyecto y quienes no
    
    \item[iii.] Obtener todos los Empleados y los Proyectos, así como los Empleados sin Proyecto y los Proyectos sin Empleado
    
    \item[iv.] Obtener los Empleados sin Proyecto y los Proyectos sin Empleados
    
    \item[v.] Obtener los Empleados que no tienen Proyecto asignado
    
    \item[vi.] Obtener los Proyectos que no tienen Empleado asignado
\end{enumerate}


\end{document}





\end{document}
